\section{Study Design}
\label{sec:study_design}

To systematically investigate the dynamics of interactive, conversational preference elicitation, we formalize our experiment using a multi-factor simulation environment. 
This design maps user interaction types against varied user objectives to evaluate the structural bounds of our system.

\subsection{Factor Structure}
Our experimental design consists of independent variables, observed dependent variables, and tightly controlled constants.

\begin{itemize}
    \item \textbf{Independent Variables (Manipulated):} 
    The primary manipulated factor is the \textit{question type} prompt-engineered into the Qwen3.6-35B-A3B interviewing agent. 
    We test three distinct interactive formats: 
    \textit{yes/no} questions, \textit{multiple-choice} questions, and \textit{open-ended} questions. 
    Additionally, the internal directive of the simulated user changes systematically on each run across 3 different personas and 6 unique goals to test the adaptive boundaries of the agent.
    \item \textbf{Dependent Variables (Observed):} 
    For each individual run, we measure conversational efficiency via the total number of \textit{turns to convergence} and the total volume of \textit{tokens exchanged}. 
    Question clarity is tracked via the \textit{user response token length}. 
    Subjective accuracy is captured via a multidimensional \textit{faithfulness score} computed using an LLM-as-judge.
    \item \textbf{Constants:} 
    To guarantee that observed changes are driven solely by the independent variables, several parameters remain fixed across all conditions: 
    the underlying corpus (a localized subset of 500 questions extracted from the StackOverflow full dataset), the core dialogue LLM (Qwen3.6-35B-A3B), the embedding encoder (Qwen3-Embedding 8B), the initial cluster count derived via K-means (used to enrich the system prompt of the interview agent), and the base behavioral persona of the simulated user.
\end{itemize}

\subsection{Outcomes and Hypothesis Testing}
We categorize our experimental objectives into primary and secondary classifications to maximize statistical transparency.
\begin{itemize}
    \item \textbf{Primary Outcome:} 
    \textit{Turns to convergence} serves as our primary outcome metric. 
    It provides an objective, direct evaluation of dialogue efficiency and establishes the headline criteria for determining which question paradigm minimizes human friction.
    \item \textbf{Secondary Outcomes:} 
    Supporting metrics include the total tokens exchanged, user response tokens, and the mean faithfulness score (the composite average of coherence, alignment, and separation metrics). 
    These support or clarify structural anomalies observed in the primary outcome but do not dictate the headline claim.
\end{itemize}

We formally evaluate the following hypothesis: 
\begin{quote}
\textit{Multiple-choice questions represent the optimal operational tradeoff between user cognitive load (response complexity) and information gained per turn, leading to superior efficiency compared to yes/no or open-ended configurations.}
\end{quote}

\subsection{Experimental Matrix}
To test this hypothesis, we execute $N=108$ distinct experimental runs. 
The three question types are balanced uniformly, resulting in 36 independent trials per interactive condition. 
The simulation environment coordinates a matrix of 18 unique variations, blending 3 distinct user personas with 6 target goals.
Each specific combination is executed exactly 2 times to ensure statistical reliability while mitigating runtime noise.

\subsection{Decision Rule}
To establish statistical significance without relying on rigid parametric assumptions, we employ a non-parametric bootstrapping strategy. 
For each pairwise comparison between interactive question conditions (e.g., Multiple-Choice vs. Open-Ended), we calculate the absolute difference in sample means and generate a $95\%$ bootstrap confidence interval using $10000$ resamples.
We conclude that Condition A is significantly more efficient than Condition B if and only if the $95\%$ confidence interval for the mean difference $(A - B)$ strictly excludes zero. 

\subsection{Baseline and Experimental Limitations}
The baseline condition embeds documents using the Qwen3-Embedding model under standard uninstructed configurations, skipping the interview stage and omitting task-specific metadata from the encoder prompt. 

This introduces an architectural confound noted explicitly in our protocol: 
the baseline differs from experimental conditions by both the \textit{absence of user preference elicitation} and the \textit{absence of task instructions during embedding generation}. 
As a result, our study design cannot mathematically isolate whether an improvement in downstream cluster faithfulness is uniquely driven by dialogue elicitation, instruction-guided embedding mechanics, or their combined synergy. 
Isolating these independent effects remains an objective for future iterations.

\subsection{Simplifying Operational Assumptions}
\label{subsec:assumptions}
To maintain a mathematically tractable and structured evaluation framework, the experimental design operates under four foundational operational assumptions:
\begin{itemize}
    \item \textbf{User Responsiveness:} 
    Simulated users are assumed to be fully responsive. When prompted, the agent always provides an explicit preference or definitively assigns a datapoint to a cluster. 
    Conversational abstentions, ambiguous inputs, or ``I do not know'' responses are explicitly out of scope.
    \item \textbf{Preference Consistency:} 
    User preferences are modeled as perfectly stable within a single conversational session. 
    The user agent will not contradict its earlier responses or exhibit shifting criteria mid-loop.
    \item \textbf{Modality Restriction:} 
    The system architecture operates strictly on a text-only data matrix. 
    Multimodal document processing and heterogeneous data inputs are considered out of scope.
    \item \textbf{Offline Clustering Topology:} 
    The framework calculates a single, absolute clustering configuration at the conclusion of the dialogue loop. 
    Partitioning is not executed incrementally during active conversation.
\end{itemize}